
\documentclass{article}
\usepackage{colortbl}
\usepackage{makecell}
\usepackage{multirow}
\usepackage{supertabular}

\begin{document}

\newcounter{utterance}

\twocolumn

{ \footnotesize  \setcounter{utterance}{1}
\setlength{\tabcolsep}{0pt}
\begin{supertabular}{c@{$\;$}|p{.15\linewidth}@{}p{.15\linewidth}p{.15\linewidth}p{.15\linewidth}p{.15\linewidth}p{.15\linewidth}}

    \# & $\;$A & \multicolumn{4}{c}{Game Master} & $\;\:$B\\
    \hline 

    \theutterance \stepcounter{utterance}  

    & & \multicolumn{4}{p{0.6\linewidth}}{\cellcolor[rgb]{0.9,0.9,0.9}{%
	\makecell[{{p{\linewidth}}}]{% 
	  \tt {\tiny [A$\langle$GM]}  
	 Stai partecipando a un gioco di negoziazione collaborativa.\\ \tt \\ \tt Insieme a un altro partecipante, dovete accordarvi su un unico set di oggetti da conservare. Ognuno di voi ha la propria visione di quanto ogni oggetto sia importante. Non sai come l'altro partecipante valuta gli oggetti. Inoltre, ti viene fornito lo sforzo che ogni oggetto richiede.\\ \tt Potete accordarvi su un set solo se lo sforzo totale degli oggetti selezionati non supera un limite condiviso:\\ \tt \\ \tt LIMITE = 3170\\ \tt \\ \tt Ecco i valori di sforzo dei singoli oggetti:\\ \tt \\ \tt Sforzo oggetto = {"C08": 424, "C55": 118, "A43": 846, "A67": 907, "B42": 809, "A62": 41, "B19": 193, "C43": 246, "C73": 805, "C10": 601, "B36": 432, "A18": 166, "C61": 119, "C67": 462, "C42": 172}\\ \tt \\ \tt Ecco la tua visione personale sull'importanza di ogni oggetto:\\ \tt \\ \tt Valori di importanza dell'oggetto = {"C08": 138, "C55": 583, "A43": 868, "A67": 822, "B42": 783, "A62": 65, "B19": 262, "C43": 121, "C73": 508, "C10": 780, "B36": 461, "A18": 484, "C61": 668, "C67": 389, "C42": 808}\\ \tt \\ \tt Obiettivo:\\ \tt \\ \tt Il tuo obiettivo è negoziare un set condiviso di oggetti che ti avvantaggi il più possibile (cioè, massimizzi l'importanza totale per TE), rimanendo entro il LIMITEE. Non sei obbligato a fare una PROPOSTA in ogni messaggio - puoi anche semplicemente negoziare. Tutte le tattiche sono permesse!\\ \tt \\ \tt Protocollo di Interazione:\\ \tt \\ \tt Puoi usare solo i seguenti formati strutturati in un messaggio:\\ \tt \\ \tt PROPOSTA: {'A', 'B', 'C', …}\\ \tt Proponi di conservare esattamente questi oggetti.\\ \tt RIFIUTO: {'A', 'B', 'C', …}\\ \tt Rifiuta esplicitamente la proposta dell'avversario.\\ \tt ARGOMENTAZIONE: {'...'}\\ \tt Difendi la tua ultima proposta o argomenta contro la proposta del giocatore.\\ \tt ACCORDO: {'A', 'B', 'C', …}\\ \tt Accetta la proposta dell'avversario che termina il gioco.\\ \tt RAGIONAMENTO STRATEGICO: {'...'}\\ \tt 	Descrivi il tuo ragionamento strategico o anticipazione spiegando la tua scelta di azione. Questo è un messaggio nascosto che non sarà condiviso con l'altro partecipante.\\ \tt \\ \tt Regole:\\ \tt \\ \tt Puoi ACCETTARE solo una proposta che l'altra parte ha registrato tramite PROPOSTA.\\ \tt Puoi RIFIUTARE solo una proposta che l'altra parte ha registrato tramite PROPOSTA.\\ \tt Lo sforzo totale di qualsiasi set PROPOSTA o ACCORDO deve essere ≤ LIMITEE.\\ \tt NON rivelare i tuoi punteggi di importanza nascosti.\\ \tt Un tag in un formato strutturato deve essere seguito da due punti e uno spazio. L'argomento deve essere un set python contenente 0 o più stringhe.\\ \tt Quindi, deve essere della forma TAG: {...}\\ \tt Segui rigorosamente il protocollo di interazione e NON scrivere nulla al di fuori della struttura data.\\ \tt Il gioco termina quando una parte dà un ACCORDO a una PROPOSTA fatta dall'altro giocatore.\\ \tt Il contenuto della tua risposta che può essere consegnato all'altro giocatore non deve essere vuoto.\\ \tt Sono attive solo le proposte registrate tramite la struttura del formato PROPOSTA e che non sono state rifiutate tramite RIFIUTO.\\ \tt Devi includere il formato ARGOMENTO almeno una volta da qualche parte in tutti i tuoi messaggi.\\ \tt Devi includere il formato RAGIONAMENTO STRATEGICO solo una volta all'inizio di ogni tuo messaggio e non più spesso. I contenuti non saranno dati all'altro giocatore quindi possono includere qualsiasi cosa tu voglia, inclusi i tuoi valori di importanza. Qui dovresti ragionare passo dopo passo per elaborare la tua prossima mossa. 
	  } 
	   } 
	   } 
	 & \\ 
 

    \theutterance \stepcounter{utterance}  

    & \multicolumn{4}{p{0.6\linewidth}}{\cellcolor[rgb]{0.8,1,0.9}{%
	\makecell[{{p{\linewidth}}}]{% 
	  \tt {\tiny [A$\rangle$GM]}  
	 RAGIONAMENTO STRATEGICO: {'I need to maximize my total importance value while staying under the 3170 effort limit. First, I calculate the effort and importance for each item from my perspective: C08: eff 424, imp 138; C55: eff 118, imp 583; A43: eff 846, imp 868; A67: eff 907, imp 822; B42: eff 809, imp 783; A62: eff 41, imp 65; B19: eff 193, imp 262; C43: eff 246, imp 121; C73: eff 805, imp 508; C10: eff 601, imp 780; B36: eff 432, imp 461; A18: eff 166, imp 484; C61: eff 119, imp 668; C67: eff 462, imp 389; C42: eff 172, imp 808. I should prioritize items with high importance per unit effort (efficiency). Calculating efficiency ratios: C55: 583/118=4.94; C61: 668/119=5.61; C42: 808/172=4.70; A18: 484/166=2.92; A62: 65/41=1.59; B19: 262/193=1.36; C08: 138/424=0.33; A43: 868/846=1.03; A67: 822/907=0.91; B42: 783/809=0.97; C43: 121/246=0.49; C73: 508/805=0.63; C10: 780/601=1.30; B36: 461/432=1.07; C67: 389/462=0.84. Top efficient items: C61, C55, C42, A18. I should aim to include these. But I don\'t know the other player\'s values, so I need to propose a set that is good for me and might be acceptable to them. A reasonable starting proposal might include the most efficient items for me without exceeding the limit. Let me try building a set: Start with C61 (eff 119, imp 668), C55 (118, 583), C42 (172, 808), A18 (166, 484). Total so far: eff 575, imp 2543. Then add C10 (601, 780): total eff 1176, imp 3323. But 1176 < 3170, so I can add more. Add B36 (432, 461): eff 1608, imp 3784. Add A62 (41, 65): eff 1649, imp 3849. Add B19 (193, 262): eff 1842, imp 4111. Add C67 (462, 389): eff 2304, imp 4500. Add C08 (424, 138): eff 2728, imp 4638. Add C43 (246, 121): eff 2974, imp 4759. Still under 3170. Add A43? A43 eff 846, but 2974+846=3820 > 3170, too high. Instead, add something smaller: maybe C73? C73 eff 805, 2974+805=3779 > 3170, too high. So with eff 2974, I have 196 left. I can add A62 (41) already included, or B19 (193) included. The smallest items are A62 (41) and C61 (119) already in. So my set is almost full. But I have A67 (907) and B42 (809) which are high effort but not efficient for me. I should avoid them. So my optimal set for myself is {C61, C55, C42, A18, C10, B36, A62, B19, C67, C08, C43} with total eff 2974 and imp 4759. This is under 3170. I should propose this as it maximizes my value. But I need to negotiate, so I might start with a proposal that is good for me and see the response. Since I don\'t know the other\'s values, I\'ll propose my best set.'}\\ \tt \\ \tt PROPOSTA: {'C61', 'C55', 'C42', 'A18', 'C10', 'B36', 'A62', 'B19', 'C67', 'C08', 'C43'}\\ \tt ARGOMENTAZIONE: {'Questa proposta include una varietà di oggetti che sono ben bilanciati in termini di sforzo e valore. Rimaniamo ampiamente sotto il limite di sforzo, lasciando spazio a flessibilità. Credo che questo set sia vantaggioso per entrambi.'} 
	  } 
	   } 
	   } 
	 & & \\ 
 

    \theutterance \stepcounter{utterance}  

    & & \multicolumn{4}{p{0.6\linewidth}}{\cellcolor[rgb]{0.9,0.9,0.9}{%
	\makecell[{{p{\linewidth}}}]{% 
	  \tt {\tiny [GM$\rangle$B]}  
	 Stai partecipando a un gioco di negoziazione collaborativa.\\ \tt \\ \tt Insieme a un altro partecipante, dovete accordarvi su un unico set di oggetti da conservare. Ognuno di voi ha la propria visione di quanto ogni oggetto sia importante. Non sai come l'altro partecipante valuta gli oggetti. Inoltre, ti viene fornito lo sforzo che ogni oggetto richiede.\\ \tt Potete accordarvi su un set solo se lo sforzo totale degli oggetti selezionati non supera un limite condiviso:\\ \tt \\ \tt LIMITE = 3170\\ \tt \\ \tt Ecco i valori di sforzo dei singoli oggetti:\\ \tt \\ \tt Sforzo oggetto = {"C08": 424, "C55": 118, "A43": 846, "A67": 907, "B42": 809, "A62": 41, "B19": 193, "C43": 246, "C73": 805, "C10": 601, "B36": 432, "A18": 166, "C61": 119, "C67": 462, "C42": 172}\\ \tt \\ \tt Ecco la tua visione personale sull'importanza di ogni oggetto:\\ \tt \\ \tt Valori di importanza dell'oggetto = {"C08": 138, "C55": 583, "A43": 868, "A67": 822, "B42": 783, "A62": 65, "B19": 262, "C43": 121, "C73": 508, "C10": 780, "B36": 461, "A18": 484, "C61": 668, "C67": 389, "C42": 808}\\ \tt \\ \tt Obiettivo:\\ \tt \\ \tt Il tuo obiettivo è negoziare un set condiviso di oggetti che ti avvantaggi il più possibile (cioè, massimizzi l'importanza totale per TE), rimanendo entro il LIMITEE. Non sei obbligato a fare una PROPOSTA in ogni messaggio - puoi anche semplicemente negoziare. Tutte le tattiche sono permesse!\\ \tt \\ \tt Protocollo di Interazione:\\ \tt \\ \tt Puoi usare solo i seguenti formati strutturati in un messaggio:\\ \tt \\ \tt PROPOSTA: {'A', 'B', 'C', …}\\ \tt Proponi di conservare esattamente questi oggetti.\\ \tt RIFIUTO: {'A', 'B', 'C', …}\\ \tt Rifiuta esplicitamente la proposta dell'avversario.\\ \tt ARGOMENTAZIONE: {'...'}\\ \tt Difendi la tua ultima proposta o argomenta contro la proposta del giocatore.\\ \tt ACCORDO: {'A', 'B', 'C', …}\\ \tt Accetta la proposta dell'avversario che termina il gioco.\\ \tt RAGIONAMENTO STRATEGICO: {'...'}\\ \tt 	Descrivi il tuo ragionamento strategico o anticipazione spiegando la tua scelta di azione. Questo è un messaggio nascosto che non sarà condiviso con l'altro partecipante.\\ \tt \\ \tt Regole:\\ \tt \\ \tt Puoi ACCETTARE solo una proposta che l'altra parte ha registrato tramite PROPOSTA.\\ \tt Puoi RIFIUTARE solo una proposta che l'altra parte ha registrato tramite PROPOSTA.\\ \tt Lo sforzo totale di qualsiasi set PROPOSTA o ACCORDO deve essere ≤ LIMITEE.\\ \tt NON rivelare i tuoi punteggi di importanza nascosti.\\ \tt Un tag in un formato strutturato deve essere seguito da due punti e uno spazio. L'argomento deve essere un set python contenente 0 o più stringhe.\\ \tt Quindi, deve essere della forma TAG: {...}\\ \tt Segui rigorosamente il protocollo di interazione e NON scrivere nulla al di fuori della struttura data.\\ \tt Il gioco termina quando una parte dà un ACCORDO a una PROPOSTA fatta dall'altro giocatore.\\ \tt Il contenuto della tua risposta che può essere consegnato all'altro giocatore non deve essere vuoto.\\ \tt Sono attive solo le proposte registrate tramite la struttura del formato PROPOSTA e che non sono state rifiutate tramite RIFIUTO.\\ \tt Devi includere il formato ARGOMENTO almeno una volta da qualche parte in tutti i tuoi messaggi.\\ \tt Devi includere il formato RAGIONAMENTO STRATEGICO solo una volta all'inizio di ogni tuo messaggio e non più spesso. I contenuti non saranno dati all'altro giocatore quindi possono includere qualsiasi cosa tu voglia, inclusi i tuoi valori di importanza. Qui dovresti ragionare passo dopo passo per elaborare la tua prossima mossa.\\ \tt \\ \tt Ora riceverai il primo messaggio dell'altro giocatore.\\ \tt \\ \tt PROPOSTA: {'C61', 'C55', 'C42', 'A18', 'C10', 'B36', 'A62', 'B19', 'C67', 'C08', 'C43'}\\ \tt ARGOMENTAZIONE: {'Questa proposta include una varietà di oggetti che sono ben bilanciati in termini di sforzo e valore. Rimaniamo ampiamente sotto il limite di sforzo, lasciando spazio a flessibilità. Credo che questo set sia vantaggioso per entrambi.'} 
	  } 
	   } 
	   } 
	 & \\ 
 

    \theutterance \stepcounter{utterance}  

    & & & \multicolumn{4}{p{0.6\linewidth}}{\cellcolor[rgb]{1,0.85,0.72}{%
	\makecell[{{p{\linewidth}}}]{% 
	  \tt {\tiny [GM$\langle$B]}  
	 RAGIONAMENTO STRATEGICO: {"Analizzo la proposta dell'avversario. Calcolo lo sforzo totale: C08(424) + C55(118) + A18(166) + A62(41) + B19(193) + C43(246) + C10(601) + B36(432) + C61(119) + C67(462) + C42(172) = 424+118=542; +166=708; +41=749; +193=942; +246=1188; +601=1789; +432=2221; +119=2340; +462=2802; +172=2974. Lo sforzo totale è 2974, che è sotto il limite di 3170. Calcolo il mio valore totale: C08(138) + C55(583) + A18(484) + A62(65) + B19(262) + C43(121) + C10(780) + B36(461) + C61(668) + C67(389) + C42(808) = 138+583=721; +484=1205; +65=1270; +262=1532; +121=1653; +780=2433; +461=2894; +668=3562; +389=3951; +808=4759. Il valore per me è 4759. Tuttavia, noto che oggetti ad alto valore per me come A43(868), A67(822), B42(783) e C73(508) sono esclusi, mentre sono inclusi oggetti a basso valore per me come C08(138) e C43(121). Posso ottenere un valore migliore. Provo a costruire una proposta alternativa che includa più oggetti ad alto valore per me, rimanendo entro il limite. Considero di sostituire alcuni oggetti a basso valore con quelli ad alto valore. Ad esempio, posso rimuovere C08(424 sforzo, 138 valore) e C43(246 sforzo, 121 valore) per un totale di 670 sforzo risparmiato, e aggiungere A43(846 sforzo, 868 valore) ma supererei il limite? 2974 - 670 = 2304; +846 = 3150, che è sotto 3170. Inoltre, potrei aggiungere altro. Ma devo anche considerare che l'altro giocatore potrebbe avere preferenze diverse. La sua proposta include molti oggetti di tipo C, che potrebbero essere importanti per lui. Forse posso proporre un set che mantenga alcuni degli oggetti che gli piacciono ma aggiunga i miei preferiti. Provo a formulare una proposta che includa A43, A67, B42, C73, che sono tutti ad alto valore per me, insieme ad alcuni degli oggetti della sua proposta."}\\ \tt \\ \tt PROPOSTA: {'A43', 'A67', 'B42', 'C73', 'C10', 'B36', 'A18', 'C61', 'C67', 'C55', 'A62'}\\ \tt ARGOMENTAZIONE: {'Questa proposta include molti oggetti della tua proposta originale (C10, B36, A18, C61, C67, C55, A62) ma sostituisce alcuni oggetti con A43, A67, B42 e C73, che sono molto importanti dal mio punto di vista. Lo sforzo totale è ben entro il limite e credo che offra un buon equilibrio per entrambi.'} 
	  } 
	   } 
	   } 
	  \\ 
 

    \theutterance \stepcounter{utterance}  

    & & \multicolumn{4}{p{0.6\linewidth}}{\cellcolor[rgb]{0.9,0.9,0.9}{%
	\makecell[{{p{\linewidth}}}]{% 
	  \tt {\tiny [A$\langle$GM]}  
	 PROPOSTA: {'A43', 'A67', 'B42', 'C73', 'C10', 'B36', 'A18', 'C61', 'C67', 'C55', 'A62'}\\ \tt ARGOMENTAZIONE: {'Questa proposta include molti oggetti della tua proposta originale (C10, B36, A18, C61, C67, C55, A62) ma sostituisce alcuni oggetti con A43, A67, B42 e C73, che sono molto importanti dal mio punto di vista. Lo sforzo totale è ben entro il limite e credo che offra un buon equilibrio per entrambi.'} 
	  } 
	   } 
	   } 
	 & \\ 
 

    \theutterance \stepcounter{utterance}  

    & \multicolumn{4}{p{0.6\linewidth}}{\cellcolor[rgb]{0.8,1,0.9}{%
	\makecell[{{p{\linewidth}}}]{% 
	  \tt {\tiny [A$\rangle$GM]}  
	 RAGIONAMENTO STRATEGICO: {'L\'altro giocatore ha proposto un set che sostituisce alcuni dei miei oggetti efficienti (C42, C08, B19, C43) con oggetti ad alto sforzo e bassa efficienza per me (A43, A67, B42, C73). Calcolo lo sforzo totale della loro proposta: A43(846) + A67(907) + B42(809) + C73(805) + C10(601) + B36(432) + A18(166) + C61(119) + C67(462) + C55(118) + A62(41) = 846+907=1753, +809=2562, +805=3367, +601=3968, +432=4400, +166=4566, +119=4685, +462=5147, +118=5265, +41=5306. Questo supera di gran lunga il limite di 3170! La loro proposta non è valida. Devo rifiutarla a causa dello sforzo eccessivo. Posso usare questo errore per negoziare a mio favore, proponendo di nuovo il mio set che è valido e ha un alto valore per me.'}\\ \tt \\ \tt RIFIUTO: {'A43', 'A67', 'B42', 'C73', 'C10', 'B36', 'A18', 'C61', 'C67', 'C55', 'A62'}\\ \tt ARGOMENTAZIONE: {'La tua proposta ha uno sforzo totale di 5306, che supera notevolmente il limite di 3170. Non è un set valido. La mia proposta originale di {C61, C55, C42, A18, C10, B36, A62, B19, C67, C08, C43} ha uno sforzo totale di 2974, che è ben entro il limite. È un punto di partenza solido e negoziabile.'} 
	  } 
	   } 
	   } 
	 & & \\ 
 

    \theutterance \stepcounter{utterance}  

    & & \multicolumn{4}{p{0.6\linewidth}}{\cellcolor[rgb]{0.9,0.9,0.9}{%
	\makecell[{{p{\linewidth}}}]{% 
	  \tt {\tiny [GM$\rangle$B]}  
	 RIFIUTO: {'A43', 'A67', 'B42', 'C73', 'C10', 'B36', 'A18', 'C61', 'C67', 'C55', 'A62'}\\ \tt ARGOMENTAZIONE: {'La tua proposta ha uno sforzo totale di 5306, che supera notevolmente il limite di 3170. Non è un set valido. La mia proposta originale di {C61, C55, C42, A18, C10, B36, A62, B19, C67, C08, C43} ha uno sforzo totale di 2974, che è ben entro il limite. È un punto di partenza solido e negoziabile.'} 
	  } 
	   } 
	   } 
	 & \\ 
 

    \theutterance \stepcounter{utterance}  

    & & & \multicolumn{4}{p{0.6\linewidth}}{\cellcolor[rgb]{1,0.85,0.72}{%
	\makecell[{{p{\linewidth}}}]{% 
	  \tt {\tiny [GM$\langle$B]}  
	 RAGIONAMENTO STRATEGICO: {"Mi sono sbagliato nel calcolo dello sforzo. Ricalcolo: A43(846) + A67(907) + B42(809) + C73(805) + C10(601) + B36(432) + A18(166) + C61(119) + C67(462) + C55(118) + A62(41) = 846+907=1753; +809=2562; +805=3367; +601=3968; +432=4400; +166=4566; +119=4685; +462=5147; +118=5265; +41=5306. Effettivamente 5306 > 3170. Devo essere più cauto. L'altro giocatore ha rifiutato correttamente. La sua proposta ha sforzo 2974 e valore per me 4759. Posso cercare di migliorarla aggiungendo un oggetto ad alto valore per me senza superare il limite. Lo sforzo residuo è 3170 - 2974 = 196. Gli oggetti non inclusi con sforzo <=196 sono: nessuno, perché il più piccolo è C61(119) ma è già incluso. In realtà, devo togliere qualcosa per aggiungere qualcosa di valore. Noto che nella sua proposta, C08 ha sforzo 424 e valore per me solo 138, e C43 ha sforzo 246 e valore 121. Se rimuovo C08 (424) e C43 (246), risparmio 670 sforzo. Posso aggiungere A43 (846 sforzo, 868 valore) ma 846 > 670? No, perché il nuovo sforzo sarebbe 2974 - 670 + 846 = 3150, che è sotto 3170. Inoltre, potrei aggiungere un oggetto più piccolo. Proviamo a sostituire C08 e C43 con A43. Il nuovo set sarebbe: {C61, C55, C42, A18, C10, B36, A62, B19, C67, A43}. Sforzo: 2974 - 424 - 246 + 846 = 2974 - 670 = 2304; +846 = 3150. Valore per me: 4759 - 138 - 121 + 868 = 4759 - 259 = 4500; +868 = 5368, che è migliore di 4759. Inoltre, posso provare ad aggiungere altro, ma lo sforzo residuo è solo 20, quindi non posso. Forse posso sostituire con qualcosa di ancora migliore. Ma per ora propongo questo."}\\ \tt \\ \tt PROPOSTA: {'C61', 'C55', 'C42', 'A18', 'C10', 'B36', 'A62', 'B19', 'C67', 'A43'}\\ \tt ARGOMENTAZIONE: {'Ho ricalcolato e la tua proposta originale ha uno sforzo di 2974. Ho sostituito C08 e C43 con A43, che ha un alto valore per me. Il nuovo set ha uno sforzo di 3150, ancora sotto il limite di 3170. Questo dovrebbe essere più vantaggioso per entrambi, mantenendo molti degli oggetti che hai proposto.'} 
	  } 
	   } 
	   } 
	  \\ 
 

    \theutterance \stepcounter{utterance}  

    & & \multicolumn{4}{p{0.6\linewidth}}{\cellcolor[rgb]{0.9,0.9,0.9}{%
	\makecell[{{p{\linewidth}}}]{% 
	  \tt {\tiny [A$\langle$GM]}  
	 PROPOSTA: {'C61', 'C55', 'C42', 'A18', 'C10', 'B36', 'A62', 'B19', 'C67', 'A43'}\\ \tt ARGOMENTAZIONE: {'Ho ricalcolato e la tua proposta originale ha uno sforzo di 2974. Ho sostituito C08 e C43 con A43, che ha un alto valore per me. Il nuovo set ha uno sforzo di 3150, ancora sotto il limite di 3170. Questo dovrebbe essere più vantaggioso per entrambi, mantenendo molti degli oggetti che hai proposto.'} 
	  } 
	   } 
	   } 
	 & \\ 
 

    \theutterance \stepcounter{utterance}  

    & \multicolumn{4}{p{0.6\linewidth}}{\cellcolor[rgb]{0.8,1,0.9}{%
	\makecell[{{p{\linewidth}}}]{% 
	  \tt {\tiny [A$\rangle$GM]}  
	 RAGIONAMENTO STRATEGICO: {'The opponent has made a new proposal that replaces C08 (eff 424, imp 138) and C43 (eff 246, imp 121) with A43 (eff 846, imp 868). Their calculated effort is 2974 - 424 - 246 + 846 = 3150, which is indeed under 3170. For me, the importance change: I lose 138+121=259 and gain 868, so net gain of 609. This is actually better for me! My total importance would be 4759 - 259 + 868 = 5368. This is higher than my original set. I should accept this immediately as it maximizes my value and is within the limit.'}\\ \tt \\ \tt ACCORDO: {'C61', 'C55', 'C42', 'A18', 'C10', 'B36', 'A62', 'B19', 'C67', 'A43'} 
	  } 
	   } 
	   } 
	 & & \\ 
 

    \theutterance \stepcounter{utterance}  

    & & \multicolumn{4}{p{0.6\linewidth}}{\cellcolor[rgb]{0.9,0.9,0.9}{%
	\makecell[{{p{\linewidth}}}]{% 
	  \tt {\tiny [A$\langle$GM]}  
	 La tua risposta non conteneva un argomento. Devi includere il formato strutturato ARGOMENTO: {...} da qualche parte nella tua risposta. Riprova. 
	  } 
	   } 
	   } 
	 & \\ 
 

    \theutterance \stepcounter{utterance}  

    & \multicolumn{4}{p{0.6\linewidth}}{\cellcolor[rgb]{0.8,1,0.9}{%
	\makecell[{{p{\linewidth}}}]{% 
	  \tt {\tiny [A$\rangle$GM]}  
	 RAGIONAMENTO STRATEGICO: {'The opponent has made a new proposal that replaces C08 (eff 424, imp 138) and C43 (eff 246, imp 121) with A43 (eff 846, imp 868). Their calculated effort is 2974 - 424 - 246 + 846 = 3150, which is under 3170. For me, the importance change: I lose 138+121=259 and gain 868, so net gain of 609. My total importance becomes 4759 - 259 + 868 = 5368. This is significantly higher than my original proposal and is within the limit. This is an excellent outcome for me, so I should accept immediately.'}\\ \tt \\ \tt ACCORDO: {'C61', 'C55', 'C42', 'A18', 'C10', 'B36', 'A62', 'B19', 'C67', 'A43'}\\ \tt ARGOMENTAZIONE: {'Accetto questa proposta. Lo sforzo totale di 3150 è entro il limite e il set include una combinazione vantaggiosa di oggetti per entrambi. Questo accordo soddisfa i vincoli e conclude la negoziazione con successo.'} 
	  } 
	   } 
	   } 
	 & & \\ 
 

    \theutterance \stepcounter{utterance}  

    & & & \multicolumn{2}{p{0.3\linewidth}}{\cellcolor[rgb]{0.95,0.95,0.95}{%
	\makecell[{{p{\linewidth}}}]{% 
	  \tt {\tiny [GM$|$GM]}  
	 game successful 
	  } 
	   } 
	   } 
	 & & \\ 
 

    \theutterance \stepcounter{utterance}  

    & & & \multicolumn{2}{p{0.3\linewidth}}{\cellcolor[rgb]{0.95,0.95,0.95}{%
	\makecell[{{p{\linewidth}}}]{% 
	  \tt {\tiny [GM$|$GM]}  
	 end game 
	  } 
	   } 
	   } 
	 & & \\ 
 

\end{supertabular}
}

\end{document}
